%% kapa: exam=1 round=34 year=2023 period=2 form=A subject=accounting subject_ko=회계학 seq=17 number=57
%% source: Q-Net 감정평가사 기출문제 (cst00309) artlSeq=5211789 file=raw/1008/5211789/2023년도 제34회 감정평가사 1차시험 2교시.pdf
%% conversion: kordoc 4.24.51 -> tools/gen1.py (auto); review=proofread-sonnet; flags=none
\begin{question}{57}
금융상품에 관한 설명으로 옳지 않은 것은?
\begin{choices}
  \item 종류별로 금융상품을 공시하는 경우에는 공시하는 정보의 특성에 맞게, 금융상품의 특성을 고려하여 금융상품을 종류별로 분류하여야 한다.
  \item 기타포괄손익-공정가치로 측정하는 금융자산의 장부금액은 손실충당금에 의해 감소되지 않는다.
  \item 당기손익-공정가치로 측정되는 지분상품은 후속적 공정가치 변동을 기타포괄손익으로 표시하도록 최초 인식시점에 선택할 수 있다.
  \item 금융자산과 금융부채를 상계하면 손익이 발생할 수 있다.
  \item 금융자산의 회수를 합리적으로 예상할 수 없는 경우에는 해당 금융자산의 총 장부금액을 직접 줄인다.
\end{choices}
\end{question}
\answer{4}
